% !TEX root = ../tex/6-1.tex
% 第6章：本文は青、補足は緑
\newcommand{\asmain}{%
    \setbeamercolor{frametitle}{fg=white,bg=main}%
    \setbeamercolor{structure}{fg=main}%
    \setbeamercolor{itemize item}{fg=main}%
    \setbeamercolor{itemize subitem}{fg=main}%
    \setbeamercolor{enumerate item}{fg=main}%
}
\newcommand{\assuppl}{%
    \setbeamercolor{frametitle}{fg=white,bg=examplecol}%
    \setbeamercolor{structure}{fg=examplecol}%
    \setbeamercolor{itemize item}{fg=examplecol}%
    \setbeamercolor{itemize subitem}{fg=examplecol}%
    \setbeamercolor{enumerate item}{fg=examplecol}%
}
% フレームタイトルの「補足」だけを□で囲む
\newcommand{\supplmark}{%
    \setlength{\fboxsep}{1.5pt}%
    \setlength{\fboxrule}{0.7pt}%
    \fbox{\textbf{補足}}%
}
\asmain


%=============================================
\begin{frame}{第5章から第6章へ}
    \begin{block}{ここまで（第2--5章）}
        変分（DDPM）、スコア（DSM / Score SDE）、\\
        フロー（FM / CFM）を、それぞれ別の動機から見てきた。
    \end{block}
    \begin{block}{ここから（第6章）}
        見た目の違いの奥に、共通原理があることを示す。\\
        鍵は「周辺の扱いづらい目標」を「条件付きの扱いやすい教師」で学ぶこと。
    \end{block}
\end{frame}
%=============================================
\asmain
\begin{frame}{この章の学習目標（導入）}
    \begin{itemize}
        \item 変分／スコア／フローに共通する\textbf{条件付け原理}を認識する
        \item 4つの予測型（clean / noise / score / velocity）を理解する
        \item 訓練設計を分解する：スケジュール・予測対象・$p_{\mathrm{time}}$・重み$\omega(t)$
        \item 「同値」の意味を区別する（オラクル／目的関数／周辺経路／実務アルゴリズム）
    \end{itemize}
    \begin{block}{今日（6.1）の焦点}
        条件付けトリック：扱いにくい周辺目標を、条件付き目標で回復する。
    \end{block}
\end{frame}
%=============================================
\asmain
\begin{frame}{6.1 の位置づけ}
    \begin{block}{この節の役割}
        三視点の訓練目標が、同じ「条件付け」でtractableになることを示す
    \end{block}
    \begin{enumerate}
        \item \textbf{扱いにくい周辺目標}\\
        $J_{\mathrm{KL}}$、$J_{\mathrm{SM}}$、$J_{\mathrm{FM}}$
        \item \textbf{条件付けトリック}\\
        $J_{\mathrm{CKL}}$、$J_{\mathrm{DSM}}$、$J_{\mathrm{CFM}}$（定数差で同値）
        \item \textbf{共通パターン（式6.1.1）}\\
        周辺オラクル＝$p(\cdot|\vx_t)$上の条件付きオラクルの平均
        \item \textbf{有限データでの具体化（任意）}\\
        事後重み$w_i$と記憶化の注意
    \end{enumerate}
\end{frame}
%=============================================
\section{三視点の扱いにくい目標}
\asmain
\begin{frame}{区切り：三視点の扱いにくい目標}
    \tableofcontents[currentsection]
\end{frame}
%=============================================
\asmain
\begin{frame}{三つの起源、三つの訓練目標}
    これまで別々に見えていた三視点は、固定$t$で次の目標に至る：
    \begin{itemize}
        \item \textbf{変分}：逆遷移$p(\vx_{t-\Delta t}|\vx_t)$を$p_\phi$で近似
        \item \textbf{スコア}：周辺スコア$\nabla_{\vx}\log p_t(\vx_t)$を$s_\phi$で近似
        \item \textbf{フロー}：周辺速度$\bm{v}_t(\vx_t)$を$\bm{v}_\phi$で近似
    \end{itemize}
    \begin{alertblock}{一見すると絶望的}
        いずれも「オラクル量」を直接必要とし、一般には計算不能。
    \end{alertblock}
\end{frame}
%=============================================
\asmain
\begin{frame}{変分：$J_{\mathrm{KL}}$}
    パラメータ化密度$p_\phi(\vx_{t-\Delta t}|\vx_t)$で\\
    オラクル逆遷移$p(\vx_{t-\Delta t}|\vx_t)$を近似する：
    \[
        J_{\mathrm{KL}}(\phi)
        :=
        \mathbb{E}_{p_t(\vx_t)}
        \Bigl[
        D_{\mathrm{KL}}
        \bigl(
        p(\vx_{t-\Delta t}|\vx_t)
        \,\big\|\,
        p_\phi(\vx_{t-\Delta t}|\vx_t)
        \bigr)
        \Bigr]
    \]
    \begin{block}{困難}
        真の逆遷移$p(\vx_{t-\Delta t}|\vx_t)$は一般に未知。
    \end{block}
\end{frame}
%=============================================
\asmain
\begin{frame}{スコア：$J_{\mathrm{SM}}$}
    スコアモデル$s_\phi(\vx_t,t)$で周辺スコアを当てる：
    \[
        J_{\mathrm{SM}}(\phi)
        :=
        \mathbb{E}_{p_t(\vx_t)}
        \bigl\|
        s_\phi(\vx_t,t)
        -
        \nabla_{\vx}\log p_t(\vx_t)
        \bigr\|_{2}^{2}
    \]
    \begin{block}{困難}
        周辺密度$p_t$のスコアは閉じた形では書けない。
    \end{block}
\end{frame}
%=============================================
\asmain
\begin{frame}{フロー：$J_{\mathrm{FM}}$}
    速度モデル$\bm{v}_\phi(\vx_t,t)$でオラクル速度（例：式5.2.9）を当てる：
    \[
        J_{\mathrm{FM}}(\phi)
        :=
        \mathbb{E}_{p_t(\vx_t)}
        \bigl\|
        \bm{v}_\phi(\vx_t,t)
        -
        \bm{v}_t(\vx_t)
        \bigr\|_{2}^{2}
    \]
    \begin{block}{困難}
        周辺速度$\bm{v}_t(\vx_t)$も、一般には未知の期待値。
    \end{block}
\end{frame}
%=============================================
\asmain
\begin{frame}{共通の行き詰まり}
    \begin{columns}[T,onlytextwidth]
        \begin{column}{0.32\textwidth}
            \begin{block}{変分}
                $p(\vx_{t-\Delta t}|\vx_t)$
            \end{block}
        \end{column}
        \begin{column}{0.32\textwidth}
            \begin{block}{スコア}
                $\nabla\log p_t$
            \end{block}
        \end{column}
        \begin{column}{0.32\textwidth}
            \begin{block}{フロー}
                $\bm{v}_t(\vx_t)$
            \end{block}
        \end{column}
    \end{columns}
    \vspace{0.6em}
    \begin{alertblock}{一言で}
        「周辺の教師」が欲しいのに、周辺の教師が手に入らない。
    \end{alertblock}
    \begin{block}{次}
        データ$\vx_0$で条件付けると、教師がtractableになる。
    \end{block}
\end{frame}
%=============================================
\section{条件付けトリック}
\asmain
\begin{frame}{区切り：条件付けトリック}
    \tableofcontents[currentsection]
\end{frame}
%=============================================
\asmain
\begin{frame}{解決策：$\vx_0$で条件付ける}
    各手法が独立にたどり着いたelegantな解は同じ：
    \begin{block}{条件付けテクニック}
        未知の周辺量ではなく、既知のガウス条件付き$p_t(\vx_t|\vx_0)$の下で\\
        条件付き目標を回帰する。
    \end{block}
    \begin{itemize}
        \item 条件付き目標は閉じた形（または安い式）で書ける
        \item 周辺目標との差は、$\phi$に依存しない定数$C$
        \item したがって勾配（最小化子）は一致する
    \end{itemize}
\end{frame}
%=============================================
\asmain
\begin{frame}{変分：CKL（式2.2.3の形）}
    \[
        \resizebox{0.98\textwidth}{!}{$
        J_{\mathrm{KL}}(\phi)
        =
        \mathbb{E}_{\vx_0}
        \mathbb{E}_{p_t(\vx_t|\vx_0)}
        \Bigl[
        D_{\mathrm{KL}}
        \bigl(
        p(\vx_{t-\Delta t}|\vx_t,\vx_0)
        \,\big\|\,
        p_\phi(\vx_{t-\Delta t}|\vx_t)
        \bigr)
        \Bigr]
        +C
        $}
    \]
    右辺の$\phi$依存部を$J_{\mathrm{CKL}}(\phi)$と呼ぶ。
    \begin{itemize}
        \item 条件付き逆遷移$p(\vx_{t-\Delta t}|\vx_t,\vx_0)$はガウスで閉じる（補題2.2.2）
        \item $J_{\mathrm{KL}}=J_{\mathrm{CKL}}+C$（$C$は$\phi$に独立）
    \end{itemize}
\end{frame}
%=============================================
\asmain
\begin{frame}{スコア：DSM（式3.3.3の形）}
    \[
        \resizebox{0.98\textwidth}{!}{$
        J_{\mathrm{SM}}(\phi)
        =
        \mathbb{E}_{\vx_0}
        \mathbb{E}_{p_t(\vx_t|\vx_0)}
        \bigl\|
        s_\phi(\vx_t,t)
        -
        \nabla_{\vx_t}\log p_t(\vx_t|\vx_0)
        \bigr\|_{2}^{2}
        +C
        $}
    \]
    右辺の$\phi$依存部を$J_{\mathrm{DSM}}(\phi)$と呼ぶ。
    \begin{itemize}
        \item 条件付きスコアは$\nabla_{\vx_t}\log p_t(\vx_t|\vx_0)=-\sigma_t^{-2}(\vx_t-\alpha_t\vx_0)$
        \item 第3章のDenoising Score Matchingそのもの
    \end{itemize}
\end{frame}
%=============================================
\asmain
\begin{frame}{フロー：CFM（式5.2.10の形）}
    \[
        \resizebox{0.98\textwidth}{!}{$
        J_{\mathrm{FM}}(\phi)
        =
        \mathbb{E}_{\vx_0}
        \mathbb{E}_{p_t(\vx_t|\vx_0)}
        \bigl\|
        \bm{v}_\phi(\vx_t,t)
        -
        \bm{v}_t(\vx_t|\vx_0)
        \bigr\|_{2}^{2}
        +C
        $}
    \]
    右辺の$\phi$依存部を$J_{\mathrm{CFM}}(\phi)$と呼ぶ。
    \begin{itemize}
        \item 条件付き速度$\bm{v}_t(\vx_t|\vx_0)$はスケジュールから閉じた形
        \item 第5章のConditional Flow Matchingと同じ原理
    \end{itemize}
\end{frame}
%=============================================
\asmain
\begin{frame}{勾配同値：周辺$\leftrightarrow$条件付き}
    \renewcommand{\arraystretch}{1.2}
    \begin{tabular}{lll}
        \hline
        視点 & 周辺 & 条件付き \\
        \hline
        変分 & $J_{\mathrm{KL}}$ & $J_{\mathrm{CKL}}$ \\
        スコア & $J_{\mathrm{SM}}$ & $J_{\mathrm{DSM}}$ \\
        フロー & $J_{\mathrm{FM}}$ & $J_{\mathrm{CFM}}$ \\
        \hline
    \end{tabular}
    \vspace{0.5em}
    \begin{block}{要点}
        条件付き版は周辺版と定数だけ違う。\\
        同じモデル族の中では最小化子が一致する。
    \end{block}
    \begin{exampleblock}{実務的意味}
        訓練では条件付き教師だけを使えばよい。\\
        周辺オラクルを直接計算する必要はない。
    \end{exampleblock}
\end{frame}
%=============================================
\section{共通パターン：事後平均}
\asmain
\begin{frame}{区切り：共通パターン（事後平均）}
    \tableofcontents[currentsection]
\end{frame}
%=============================================
\asmain
\begin{frame}{理想の最小化子が回復するもの}
    モデルが十分表現力をもてば、三視点は同じ概念パターンに収束する：
    \[
        \underbrace{\text{周辺オラクル}}_{\text{一般に intractable}}
        =
        \mathbb{E}_{\vx_0\sim p(\cdot|\vx_t)}
        \Bigl[
        \underbrace{\text{条件付きオラクル}}_{\vx_0\text{が与えればtractable}}
        \Bigr]
    \]
    \begin{itemize}
        \item $\vx_0$で条件付ける$\Rightarrow$閉じた教師
        \item $\vx_t$と整合する$\vx_0$について平均$\Rightarrow$欲しい周辺オラクル
    \end{itemize}
\end{frame}
%=============================================
\asmain
\begin{frame}{式(6.1.1)：三つのオラクル恒等式}
    \[
        \begin{aligned}
        p^\star(\vx_{t-\Delta t}|\vx_t)
        &=
        \mathbb{E}_{\vx_0\sim p(\cdot|\vx_t)}
        \bigl[p(\vx_{t-\Delta t}|\vx_t,\vx_0)\bigr]
        =
        p(\vx_{t-\Delta t}|\vx_t),
        \\
        s^\star(\vx_t,t)
        &=
        \mathbb{E}_{\vx_0\sim p(\cdot|\vx_t)}
        \bigl[\nabla_{\vx_t}\log p_t(\vx_t|\vx_0)\bigr]
        =
        \nabla_{\vx_t}\log p_t(\vx_t),
        \\
        \bm{v}^\star(\vx_t,t)
        &=
        \mathbb{E}_{\vx_0\sim p(\cdot|\vx_t)}
        \bigl[\bm{v}_t(\vx_t|\vx_0)\bigr]
        =
        \bm{v}_t(\vx_t).
        \end{aligned}
        \tag{6.1.1}
    \]
\end{frame}
%=============================================
\asmain
\begin{frame}{なぜ平均になるか}
    \begin{columns}[T,onlytextwidth]
        \begin{column}{0.48\textwidth}
            \begin{block}{変分}
                条件付き逆遷移を$\vx_0$で周辺化すると\\
                真の逆遷移。\\
                KLの最小化子はその周辺。
            \end{block}
        \end{column}
        \begin{column}{0.48\textwidth}
            \begin{block}{スコア／速度}
                二乗誤差回帰の母集団最小化子は\\
                条件付き期待値。
            \end{block}
        \end{column}
    \end{columns}
    \vspace{0.5em}
    \begin{alertblock}{注意}
        これらは unrestricted な母集団オラクルの話。\\
        有限容量のニューラルネットが完全再現するとは限らない。
    \end{alertblock}
\end{frame}
%=============================================
\asmain
\begin{frame}{直観：$p(\vx_0|\vx_t)$上の平均}
    観測$\vx_t$を見たとき、あり得るきれいなデータ$\vx_0$の事後
    \[
        p(\vx_0|\vx_t)
        =
        \frac{p_t(\vx_t|\vx_0)\,p_{\mathrm{data}}(\vx_0)}
             {\int p_t(\vx_t|\tilde{\vx}_0)\,p_{\mathrm{data}}(\tilde{\vx}_0)\,d\tilde{\vx}_0}
    \]
    の上で、条件付き教師を平均している。
    \begin{block}{一言}
        拡散モデルの学習目標は、\\
        「ノイズ観測から見た、きれいなデータの事後平均」である。
    \end{block}
\end{frame}
%=============================================
\section{有限データでの具体化}
\assuppl
\begin{frame}{区切り：\supplmark\; 有限データでの具体化}
    \tableofcontents[currentsection]
\end{frame}
%=============================================
\assuppl
\begin{frame}{\supplmark：経験測度と事後重み}
    データを経験測度
    $\hat{p}_0=\frac{1}{N}\sum_{i=1}^{N}\delta_{\vx_0^{(i)}}$
    で置くと、事後は訓練点だけに乗る：
    \[
        p(\vx_0=\vx_0^{(i)}|\vx_t)
        =
        \frac{p_t(\vx_t|\vx_0^{(i)})}
             {\sum_{j=1}^{N}p_t(\vx_t|\vx_0^{(j)})}
        =:
        w_i(\vx_t,t),
        \qquad
        \sum_{i=1}^{N}w_i=1.
    \]
    任意の量$h(\vx_0)$の事後期待は
    \[
        \mathbb{E}_{\vx_0\sim p(\cdot|\vx_t)}[h(\vx_0)]
        =
        \sum_{i=1}^{N}w_i(\vx_t,t)\,h(\vx_0^{(i)}).
    \]
\end{frame}
%=============================================
\assuppl
\begin{frame}{\supplmark：変分オラクル（有限データ）}
    \[
        p^\star(\vx_{t-\Delta t}|\vx_t)
        =
        \sum_{i=1}^{N}
        w_i(\vx_t,t)\,
        p(\vx_{t-\Delta t}|\vx_t,\vx_0^{(i)})
    \]
    \begin{itemize}
        \item 各$p(\vx_{t-\Delta t}|\vx_t,\vx_0^{(i)})$はガウスブリッジで閉じる
        \item 真の逆遷移は、明示的ガウスの事後重み付き混合
    \end{itemize}
\end{frame}
%=============================================
\assuppl
\begin{frame}{\supplmark：スコアオラクル（有限データ）}
    条件付きスコア
    $\nabla_{\vx_t}\log p_t(\vx_t|\vx_0^{(i)})
    =-\sigma_t^{-2}(\vx_t-\alpha_t\vx_0^{(i)})$
    より
    \[
        s^\star(\vx_t,t)
        =
        -\frac{1}{\sigma_t^{2}}
        \Biggl(
        \vx_t
        -
        \alpha_t
        \sum_{i=1}^{N}w_i(\vx_t,t)\,\vx_0^{(i)}
        \Biggr).
    \]
    \begin{exampleblock}{読み方}
        周辺スコアは、ノイズ点から「事後平均されたきれいなデータ」へ向かう。
    \end{exampleblock}
\end{frame}
%=============================================
\assuppl
\begin{frame}{\supplmark：速度オラクル（有限データ）}
    条件付き速度
    \[
        \bm{v}_t(\vx_t|\vx_0^{(i)})
        =
        \alpha_t'\vx_0^{(i)}
        +
        \frac{\sigma_t'}{\sigma_t}
        \bigl(\vx_t-\alpha_t\vx_0^{(i)}\bigr)
    \]
    を重み付き平均すると
    \[
        \bm{v}^\star(\vx_t,t)
        =
        \sum_{i=1}^{N}
        w_i(\vx_t,t)\,
        \bm{v}_t(\vx_t|\vx_0^{(i)}).
    \]
    三視点とも、閉じた条件付き量の事後重み付き平均に帰着する。
\end{frame}
%=============================================
\assuppl
\begin{frame}{\supplmark：記憶化（memorization）への注意}
    \begin{itemize}
        \item 上式は人口分布$p_{\mathrm{data}}$ではなく、経験測度$\hat{p}_0$のオラクル
        \item 経験オラクルを全ノイズレベルで完全学習し、逆過程も誤差なしなら\\
        生成分布は$\hat{p}_0$（＝訓練点のみ）になる
        \item 理想化された意味では「完全フィット＝記憶化」
    \end{itemize}
    \begin{alertblock}{しかし実務では}
        有限容量・有限最適化・数値サンプリング誤差がある。\\
        有限データ計算は記憶化が厳密解になり得ることを示すが、\\
        実モデルが必ず記憶するとは言わない。
    \end{alertblock}
\end{frame}
%=============================================
\section{まとめ}
\asmain
\begin{frame}{区切り：6.1のまとめ}
    \tableofcontents[currentsection]
\end{frame}
%=============================================
\asmain
\begin{frame}{条件付けは偶然ではない}
    共通の条件付き書き換えは偶然ではない：
    \begin{enumerate}
        \item 訓練をtractableにする
        \item 三視点が同じ原理の異なる面だと示す
        \item 変分拡散・スコアSDE・Flow Matchingを一本の物語に繋ぐ
    \end{enumerate}
    \begin{block}{標語}
        Three perspectives, one insight.
    \end{block}
\end{frame}
%=============================================
\asmain
\begin{frame}{6.1 全体のまとめ}
    \begin{enumerate}
        \item \textbf{行き詰まり}
        $J_{\mathrm{KL}}$／$J_{\mathrm{SM}}$／$J_{\mathrm{FM}}$は周辺オラクルを要し intractable
        \item \textbf{条件付け}
        $\vx_0$条件で$J_{\mathrm{CKL}}$／$J_{\mathrm{DSM}}$／$J_{\mathrm{CFM}}$へ。定数差で同値
        \item \textbf{式(6.1.1)}
        周辺オラクル＝$p(\vx_0|\vx_t)$上の条件付きオラクル平均
        \item \textbf{有限データ}
        事後重み$w_i$の明示和；記憶化は理想解の注意書き
    \end{enumerate}
\end{frame}
%=============================================
\asmain
\begin{frame}{6.1 全体のまとめ：一本の流れ}
    \begin{block}{今日の道筋}
        扱いにくい周辺目標
        $\;\to\;$
        $\vx_0$条件付け
        $\;\to\;$
        定数差の同値
        $\;\to\;$
        事後平均としてのオラクル
    \end{block}
    \begin{block}{核心}
        欲しいのは周辺の教師。\\
        閉じた条件付き教師を回帰すれば、平均的に周辺が得られる。
    \end{block}
    \begin{block}{次へ（6.2）}
        予測型を4つに広げ、訓練損失の設計要素を分解する。
    \end{block}
\end{frame}
